\section{Dataset}
\label{sec:dataset}

\subsection{UIT-ViHSD}
We use \emph{UIT-ViHSD} \cite{luu2021vihsd}, a human-annotated corpus of Vietnamese
social-media comments, each carrying exactly one label: \textsc{clean},
\textsc{offensive} (insulting but not hate speech) or \textsc{hate} (targeting a
person or protected group).

We follow the official train/dev/test partition.  Preprocessing is deliberately
light and label-preserving: strip leading/trailing whitespace, collapse runs of
whitespace, drop control characters, and drop only rows with no usable text (2
rows in train).  We \emph{keep} duplicates and do not lowercase, strip diacritics,
normalise teencode or truncate emoji, because those surface cues are part of the
signal.  This yields 33{,}398 comments (24{,}046 / 2{,}672 / 6{,}680 for
train/dev/test); the corpus has 33{,}400 rows---the difference is exactly the 2
empty-text train rows (Table~\ref{tab:splits}).  Comments are short and uneven
(test mean 47.5 characters, median 33; \textsc{hate} about twice as long as the
others).  The annotation study reports only moderate inter-annotator agreement
($\kappa\approx0.5$), which sets a ceiling on accuracy: some test comments are
genuinely ambiguous between \textsc{offensive} and \textsc{hate}, and no recipe
removes that label noise.

\subsection{Label Mapping and Imbalance}
\label{sec:imbalance}
The native ViHSD identifiers are (0 = \textsc{clean}, 1 = \textsc{offensive},
2 = \textsc{hate}); because the minority classes are the object of study, we freeze
a canonical mapping with the rare class first ($\textsc{hate}=0$,
$\textsc{offensive}=1$, $\textsc{clean}=2$), applied once and never changed.

The distribution is skewed and stable across splits: \textsc{clean} is 82--83\% of
every split, \textsc{hate} about 10\%, and \textsc{offensive} the smallest class at
6.6--7.9\% (test: 688 \textsc{hate}, 444 \textsc{offensive}, 5{,}548
\textsc{clean}).  Two consequences follow.  First, accuracy is misleading---a
degenerate \textsc{clean}-always classifier would score about $0.83$ yet be
useless---so macro-F1 is primary.  Second, offensive and hate speech share most
surface markers, so the interesting comparison is the trade-off between
minority-class F1 and the easy \textsc{clean} class.

\subsection{Fixed Split and Reproducibility}
The split must be identical for every system, so we store it as CSV files and record
each file's SHA-256 digest in a manifest copied into every run directory; all six
systems read the same three files and label mapping, leaving the recipe as the only
varying factor, with the test split scored once per system after hyper-parameters
are frozen on validation (Table~\ref{tab:splits}).

% Fixed split used by every system; digests are the first 12 hex digits of the
% full SHA-256 recorded in data/splits/split_manifest.json.
\begin{table}[t]
\centering
\small
\setlength{\tabcolsep}{4pt}
\caption{The fixed split stored once and reread by every system.  Counts are
comments after cleaning (Sect.~\ref{sec:dataset}).  The SHA-256 digest (first
12 hex digits of the full value in the released manifest) is copied into each run
directory and re-checked before training.}
\label{tab:splits}
\begin{tabular}{lrrrrl}
\toprule
Split & Comments & HATE & OFF. & CLEAN & SHA-256 \\
\midrule
train      & 24{,}046 & 2{,}556 & 1{,}605 & 19{,}885 & \texttt{f705ffe94d41} \\
validation & 2{,}672  & 270     & 212     & 2{,}190  & \texttt{6aa552aafb92} \\
test       & 6{,}680  & 688     & 444     & 5{,}548  & \texttt{d8f9ae7ddd8c} \\
\bottomrule
\end{tabular}
\end{table}


\subsection{Duplication Across Splits}
Because we keep duplicates, we measured their cross-split prevalence.  The official
ViHSD partition reuses comments: after case-folding, $895$ of the $6{,}680$ test
comments ($13.4\%$) have an exact duplicate in train or validation and $118$
duplicated strings carry conflicting labels, a property of the public split that
can only flatter absolute scores.  Table~\ref{tab:duplicate-free} re-scores all six
systems on the $5{,}785$ duplicate-free comments: accuracy drops by $0.3$--$0.5$
points for every system while macro-F1 and its ordering barely move (M3 still
highest at $0.6654$; only the two lowest systems swap, by $0.0003$), so the
conclusions do not rest on the overlap (Sect.~\ref{sec:threats}).

% AUTO-GENERATED by scripts/08_check_leakage.py - do not edit by hand.
\begin{table}[h]
\centering
\footnotesize
\setlength{\tabcolsep}{4pt}
\caption{Full vs.\ the 5,785 duplicate-free test comments (Sect.~\ref{sec:dataset}).  Accuracy drops for every system while macro-F1 and its ordering are essentially unchanged, so the conclusions do not rest on the overlap (Sect.~\ref{sec:threats}).}
\label{tab:duplicate-free}
\begin{tabular}{lcccc}
\toprule
 & \multicolumn{2}{c}{Accuracy} & \multicolumn{2}{c}{Macro-F1} \\
\cmidrule(lr){2-3} \cmidrule(lr){4-5}
System & Full & Dup-free & Full & Dup-free \\
\midrule
M1 & 0.8734 & 0.8681 & 0.6545 & 0.6533 \\
M2 & 0.8762 & 0.8719 & 0.6575 & 0.6576 \\
M3 & 0.8754 & 0.8721 & \textbf{0.6638} & \textbf{0.6654} \\
M4 & 0.8731 & 0.8698 & 0.6590 & 0.6601 \\
M5 & \textbf{0.8783} & \textbf{0.8748} & 0.6536 & 0.6537 \\
M6 & 0.8633 & 0.8588 & 0.6604 & 0.6606 \\
\bottomrule
\end{tabular}
\end{table}

